\section{A general tensor-count and scalar-kernel bound}
\label{sec:tensor-count-scalar-bound}

\paragraph{Status and scope.}
New root derivation HD-TC01--03. Complete author and peer checks
remain pending. The bounds concern an actual inducing block of
a finite-cover sixfold source with period rank nine. In its
accepted primitive normal-core gate, the irreducibly acting
perfect normal subgroup is a central product of quasisimple
groups, and the full inducing group acts transitively on its
simple factors. Their complex module dimensions are consequently
equal; write $a\ge2$, $t$ for the number of factors and $e=a^t$.
The scalar, outer and factor-permuting actions are all retained.
No equality of original and inducing fields is assumed.

The actual base-cover and marked-moduli arguments established in
D-BI-01 and the subsequent normal-core reductions give a sphere
base, a generating product-one inertia tuple with $r\ge12$
genuine entries, and the exact identity
\[
              D:=\sum_{i=1}^r\operatorname{rank}(T_i-I)
                           =2e+6.                       \tag{TC.1}
\]
The natural factor-permutation quotient is an actual connected
cover of degree $t$. Thus its inertia permutations $\pi_i$ satisfy
\[
          \sum_i\operatorname{ind}(\pi_i)\ge2t-2,
 \qquad\operatorname{ind}(\pi)=t-\#\operatorname{cycles}(\pi).
                                                               \tag{TC.2}
\]
These are hypotheses furnished by the actual source, not
properties imposed on an arbitrary list of matrices.

\subsection*{HD-TC01. A cyclic-average tensor permutation bound}

\begin{lemma}[HD-TC01]
Let $T$ be a finite-order tensor normalizer on a tensor product
of $t$ spaces of dimension $a\ge2$, with factor permutation $\pi$.
Then its maximum-eigenspace codimension satisfies
\[
 \frac{\nu(T)}{a^t}\ge
       \bigl(1-a^{-t/2}\bigr)\frac{\operatorname{ind}(\pi)}t.
                                                               \tag{TC.3}
\]
If $\pi\ne1$, it also satisfies
$\nu(T)\ge(a-1)a^{t-1}/2\ge a^t/4$.
\end{lemma}

\begin{proof}
Choose invariant positive Hermitian forms on the irreducible
finite core factors and take their tensor product. Its scalar
uniqueness follows from irreducibility. A finite-order normalizer
is unitary for this form: its positive similitude is a positive
real root of unity. Its tensor maps can therefore be rescaled
to unitary maps, with only an overall phase remaining.
The usual tensor trace identity gives, for every integer $j$,
\[
 |\operatorname{tr}(T^j)|\le
                       a^{\#\operatorname{cycles}(\pi^j)}.
                                                               \tag{TC.4}
\]
It is the product of the traces of the operators around the
cycles, each of absolute value at most $a$.

Let $f_j$ be the number of fixed letters of $\pi^j$.
Every nontrivial permutation cycle has index at least half
its length, so
\[
 a^{\#\operatorname{cycles}(\pi^j)-t}
       \le a^{-(t-f_j)/2}
       \le\frac{f_j}{t}
                  +\left(1-\frac{f_j}{t}\right)a^{-t/2}.
                                                               \tag{TC.5}
\]
The last inequality is the endpoint chord inequality for the
convex function $x\mapsto a^{-(t-x)/2}$ on $[0,t]$.
If $L$ is the order of $T$, it is a multiple of the order
of $\pi$. The cyclic orbit-counting identity gives
$L^{-1}\sum_{j=0}^{L-1}f_j=\#\operatorname{cycles}(\pi)$.
For every eigenvalue $z$ of $T$, the Fourier formula and
(TC.4)--(TC.5) therefore bound its multiplicity by
\[
 \frac{\dim\ker(T-zI)}{a^t}
 \le\frac1L\sum_j\frac{|\operatorname{tr}(T^j)|}{a^t}
 \le\frac{\#\operatorname{cycles}(\pi)}t
       +\frac{\operatorname{ind}(\pi)}t a^{-t/2}.
                                                               \tag{TC.6}
\]
This proves (TC.3) for every target eigenvalue and every phase.

For the second assertion choose one factor cycle of length
$l\ge2$. Conjugating its weighted tensor rotation puts its
cycle product on one factor; diagonalize that product.
The resulting tensor basis is permuted up to nonzero scalar
weights. There are $a$ constant words. Every other word orbit
has length at least two and contributes at most one vector
to a specified eigenspace. Hence the maximum eigenspace on
this cycle has dimension at most $(a^l+a)/2$.
Its support fraction is at least
$(1-a^{1-l})/2\ge(a-1)/(2a)$. Tensoring with the other
cycles cannot increase the largest normalized eigenvalue
multiplicity, by convolution of the spectral multiplicities.
This proves the second bound. It is also the weighted-word
argument already used in D-ATT-08, with all short orbits retained.
\end{proof}

\begin{theorem}[HD-TC02]
Every source in the stated transitive-core gate satisfies
\[
 t\le T(a):=\max\left\{4,\left\lceil\frac{5a}{4}\right\rceil-1\right\}.
                                                               \tag{TC.7}
\]
This applies over every original imaginary quadratic field,
to arbitrary nontrivial factor modules and represented central
quotients, without a character-ratio or self-duality assumption.
\end{theorem}

\begin{proof}
There is nothing to prove if $t\le4$. Otherwise $t\ge5$
and $e=a^t\ge32$. If nine inertia entries moved factors,
HD-TC01 would give $D\ge9e/4>2e+6$, since $e>24$.
At most eight entries move factors, leaving at least four
genuine factor-fixing entries.

At such an entry with a nonscalar tensor factor, that
factor has maximum-eigenspace codimension at least one.
Spectral convolution gives moving rank at least $e/a$.
If all tensor factors are scalar, the genuine entry is
a nonidentity scalar and has full moving rank $e$.
The same lower bound holds in that case.
Sum these four bounds and use HD-TC01 and (TC.2) on the
factor-moving entries. The result is
\[
 \begin{split}
 2+\frac6e
 &\ge\left(1-a^{-t/2}\right)\left(2-\frac2t\right)+\frac4a,\\
 \frac4a&\le\frac2t+2a^{-t/2}+6a^{-t}<\frac5t.
 \end{split}                                             \tag{TC.8}
\]
For the last strict inequality, $2^t>t^2$ and $2^t>6t$
for every integer $t\ge5$. The first follows from equality
at $t=4$ and induction using $2t^2>(t+1)^2$ for $t\ge3$;
the second follows from $32>30$ and immediate induction.
Thus $a^{-t/2}<1/t$ and $6a^{-t}<1/t$.
Consequently $t<5a/4$. Taking the largest admissible integer
proves (TC.7).
\end{proof}

\subsection*{HD-TC03. Scalars and finite bounded-factor catalogues}

\begin{theorem}[HD-TC03]
Let $J$ be the actual finite inducing image and let $S$ be
its irreducibly acting perfect normal central product.
Put $P=S/Z(S)$ and $M=|\operatorname{Aut}(P)|$.
Then the scalar kernel and the whole image satisfy
\[
 |J\cap\mathbf C^\times I|\mid eM,
 \qquad |J|\le eM^2.                                    \tag{TC.9}
\]
Suppose a fixed finite catalogue of allowed simple factors
and factor modules has degree bound $A\ge2$ and automorphism
order bound $B\ge1$ for its simple factors. Put
\[
 T=T(A),\qquad E=A^T,\qquad
 Q=E\bigl(B^T T!\bigr)^2.                               \tag{TC.10}
\]
Every actual source drawn from this catalogue has
\[
 e\le E,\quad |J|\le Q,\quad r\le2E+6,
 \quad g(C)\le1+Q(E+2)                                 \tag{TC.11}
\]
for its connected Galois cover $C$ of the sphere.
In particular its remaining finite-group, scalar and
Nielsen-component search is finite and effective.
\end{theorem}

\begin{proof}
Conjugation gives $J\to\operatorname{Aut}(P)$.
If an element induces the identity on $P$, its conjugation
on $S$ has the form $s\mapsto z_s s$ with $z_s\in Z(S)$.
Centrality makes $s\mapsto z_s$ a homomorphism. Perfectness
of $S$ makes it trivial. The element therefore centralizes
$S$ and is scalar by irreducibility. Conversely a scalar
centralizes $S$, so the projective image of $J$ embeds in
$\operatorname{Aut}(P)$ and has order dividing $M$.

There cannot be two inertia entries of full moving rank:
they and the other at least ten genuine entries would
give $D\ge2e+10>2e+6$. Every entry except possibly one
therefore has eigenvalue one. Its $M$th power is scalar
by the projective embedding, and evaluation on its fixed
vector shows that this scalar is one. Thus its determinant
has order dividing $M$. Product one gives the same
determinant bound for the possible exceptional entry.
Since the actual inertia tuple generates $J$, the entire
determinant image is contained in $\mu_M$.

For a scalar $zI\in J$, its determinant is $z^e$; hence
$z^{eM}=1$. The finite scalar subgroup is cyclic and its
order divides $eM$. Multiplying its size by the order of
the projective image proves (TC.9). No rational scalar
section of the projective normalizer is needed.

By transitivity all simple factors are isomorphic, so
$P\simeq L^t$ and
$|\operatorname{Aut}(P)|=|\operatorname{Aut}(L)|^t t!$.
Together with HD-TC02 this gives the numerical bounds
$e\le E$ and $M\le B^T T!$, hence $|J|\le Q$.
Every genuine inertia contributes at least one to (TC.1),
so $r\le2e+6\le2E+6$. The Galois Riemann--Hurwitz formula
then gives
\[
 2g(C)-2=|J|\left(-2+\sum_i(1-1/o(T_i))\right)
              \le |J|(r-2)\le Q(2E+4),                  \tag{TC.12}
\]
which proves the genus bound.

For effectiveness, enumerate finite multiplication tables
of order at most $Q$, their irreducible characters of degree
at most $E$, and tuples of length at most $2E+6$ satisfying
generation and product one. All are finite algebraic or
finite-group operations. Retain the actual original-field,
Schur-index, signature and central-core conditions rather
than infer them from the group-order bound. The accepted
terminating per-component Nielsen-rank procedure D-NR
then applies to this finite list. In particular the formerly
unbounded tensor count and scalar extension do not remain
outside that finite decision procedure. No such large finite
enumeration is claimed to have been executed here.
\end{proof}

\paragraph{Strictness and remaining boundary.}
For example the ordinary label $(96,4)$ at $n=100$ has
$a=\binom{100}{4}-\binom{100}{3}$ and satisfies the earlier
Gaussian necessary long-row and field inequalities.
Its profile with $t=2a$ was not removed by a bound on tensor
count for general remaining labels; HD-TC02 now excludes it.
This is strict removal of a necessary profile, without a
claim that an actual source existed there.
For every bounded classical-factor catalogue the result
makes the remaining curve-source decision finite. Unbounded
ordinary alternating natural degrees remain a separate gate;
no nonexistence or dominance conclusion is inferred from
finiteness alone.

These inequalities concern actual finite etale local systems
away from the branch divisor and their tensor operations in
representation lambda-rings and K-theory. They make a part of
the condition-A supplier gate effective (cote16, scan050;
SIGN\_NEXT SN-C03/C04), retaining the established schemes-over-
$\mathbf Z$, EGA/SGA/FGA and motivic framework. They produce
no condition-I positive cycle and introduce no Markman,
Severn or032 premise.
